%%%%%%%%%%%%%%%%%%%%%%%%%%%%%%%%%%%%%%%%%%%%%%%%%%%%%%%%%%%%%%%%%%%%%%%%%%%%%
%% Descr:       Vorlage für Berichte der DHBW-Karlsruhe, Ein Kapitel
%% Author:      Prof. Dr. Jürgen Vollmer, vollmer@dhbw-karlsruhe.de
%% $Id: kapitel2.tex,v 1.5 2017/10/06 14:02:51 vollmer Exp $
%%  -*- coding: utf-8 -*-
%%%%%%%%%%%%%%%%%%%%%%%%%%%%%%%%%%%%%%%%%%%%%%%%%%%%%%%%%%%%%%%%%%%%%%%%%%%

\chapter{Projektverlauf und Ergebnisse}

Wie in der Projektdefinition beschrieben, war das Ziel der Studienarbeit die Entwicklung eines Smart Mirrors mit Pulsmessung und Mood Detection. Im Folgenden wird der Projektverlauf und die erzielten Ergebnisse beschrieben.

Das Display konnte erfolgreich hinter dem Spionspiegel getestet werden. Die Bildqualität sowie die Lesbarkeit der angezeigten Informationen wurden getestet und als sehr gut bewertet. Die Anbindung der Industriekamera ist erfolgt; erste Tests der Pulsmessung und Mood Detection (über OpenCV/ML-Algorithmen) konnten starten, wobei die Resultate vielversprechend waren und grundsätzlich die gewünschte Funktionalität möglich machten.

Die Softwarearchitektur basiert auf Python und geeigneten Bibliotheken für Bildauswertung und GUI. Die Datenübertragung von Kamera zu Auswertung und schließlich zur Anzeige auf dem Display funktioniert stabil im Labortest.
